% 第 11 题四棱锥三视图：正视图（左上，等腰三角形＋内部实线/虚线/各连向顶点）、侧视图（右上，
% 等腰三角形无内部线）、俯视图（左下，平行四边形＋两条对角线）。按固定源 PDF 矢量坐标换算，
% 单位 u = 25.5 pt = 0.896 cm（原图 1 m 按 25.5 pt 绘制）；坐标以正视图底边左端为原点、y 向上。
% 尺寸标注：左侧竖直「3」尺寸线（两端外向箭头、上下端线向左伸）、底边下方水平尺寸链
% 「1+1+1」与侧视图「1」（各被数字断开、6 条竖直端线）；三段图注居中于各自图形下方。
\begin{tikzpicture}[>=stealth,x=0.896cm,y=0.896cm,every node/.style={font=\small}]
% ===== 正视图 =====
\draw (0,0)--(3,0)--(1.5,3)--cycle;
\draw (2,0)--(1.5,3);
\draw[dashed] (1,0)--(1.5,3);
% 左侧竖直尺寸线（x = -0.3），两端外向箭头，数字 3
\draw[->] (-0.3,0.55)--(-0.3,0);
\draw[->] (-0.3,2.45)--(-0.3,3);
\node at (-0.3,1.5) {\(3\)};
\draw (-0.5,3)--(-0.1,3);
\draw (-0.5,0)--(-0.1,0);
% 底边下方水平尺寸链（y = -0.3）：三段 1
\draw[->] (0.45,-0.3)--(0,-0.3);
\draw[->] (0.55,-0.3)--(1,-0.3);
\node at (0.5,-0.3) {\(1\)};
\draw[->] (1.45,-0.3)--(1,-0.3);
\draw[->] (1.55,-0.3)--(2,-0.3);
\node at (1.5,-0.3) {\(1\)};
\draw[->] (2.45,-0.3)--(2,-0.3);
\draw[->] (2.55,-0.3)--(3,-0.3);
\node at (2.5,-0.3) {\(1\)};
\draw (0,-0.5)--(0,-0.1);
\draw (1,-0.5)--(1,-0.1);
\draw (2,-0.5)--(2,-0.1);
\draw (3,-0.5)--(3,-0.1);
\node[anchor=north] at (1.5,-0.77) {正视图};
% ===== 侧视图（与正视图同基线）=====
\draw (4,0)--(5,0)--(4.5,3)--cycle;
\draw[->] (4.45,-0.3)--(4,-0.3);
\draw[->] (4.55,-0.3)--(5,-0.3);
\node at (4.5,-0.3) {\(1\)};
\draw (4,-0.5)--(4,-0.1);
\draw (5,-0.5)--(5,-0.1);
\node[anchor=north] at (4.5,-0.77) {侧视图};
% ===== 俯视图（正视图正下方）=====
\draw (0,-2.5)--(2,-2.5)--(3,-1.5)--(1,-1.5)--cycle;
\draw (0,-2.5)--(3,-1.5);
\draw (2,-2.5)--(1,-1.5);
\node[anchor=north] at (1.5,-2.77) {俯视图};
\end{tikzpicture}
